%% ============================================================
%%  Chapter 8 --- Persistence Requires Governance: Long-Term Memory with mem0
%%  Production-Grade Agentic AI: LangGraph + Python Frameworks
%% ============================================================

\chapter{Chapter 8 --- Persistence Requires Governance: Long-Term Memory with mem0}

\chapterquote{A memory that outlives the session also outlives the mistake that wrote it.}{anon}

The production failure this chapter prevents is silent persistence drift: the
condition where an agent stores trivial chatter, stale user facts, duplicated
preferences, or cross-tenant records in a long-term memory system and then
faithfully retrieves them weeks later as if they were authoritative. The cost
is not merely lower answer quality. It is compliance exposure, retrieval noise
that degrades relevance, operational bloat in the vector store, and the most
expensive failure of all: a system that can only explain why it answered badly
by admitting that it remembered the wrong thing.

The solution has four layers. The first layer is \emph{construction}: a single
mem0 client, built once from validated configuration, with explicit deployment
mode and backend choices. The second layer is \emph{governed writes}: a typed
memory taxonomy, a five-check write gate, and metadata that gives every
persisted record a lifecycle. The third layer is \emph{governed reads}: a
retrieval pipeline that reuses Chapter~7's six-step injection contract,
reranks what mem0 returns, and injects only the small set of memories worth
paying for in prompt tokens. The fourth layer is \emph{operations}: expiry,
pruning, inspection, optimistic locking, audit events, tenant isolation, and
tests that prove the policy behaves under load rather than merely compiling.

The sections follow dependency order. Section~\ref{sec:deployment-mode}
establishes the interface and deployment boundary, because the rest of the
chapter treats mem0 as a singleton dependency rather than an ad-hoc helper.
Section~\ref{sec:config-not-code} fixes configuration next, because write and
read policy are meaningless if the underlying store is unstable.
Sections~\ref{sec:type-explicit} through~\ref{sec:semantic-extraction} govern
what enters memory and what mem0 does with it after a write.
Section~\ref{sec:retrieval-narrow} shows how long-term memory reuses the
six-step pipeline that Chapter~7 introduced as the injection contract.
The remaining sections move outward into lifecycle and operability: expiry,
pruning, inspection, locking, auditing, multi-tenancy, and tests. That
ordering matters because you cannot reason about deletion until you know what
was written, and you cannot test retrieval until the write and read contracts
are settled.

This chapter introduces more than three persistent-memory terms, so
Table~\ref{tab:mem0-taxonomy} defines them before any implementation detail.

\begin{table}[ht]
\centering
\caption{Persistent-memory taxonomy introduced in Chapter~8}
\label{tab:mem0-taxonomy}
\begin{tabular}{llll}
\toprule
\textbf{Term} & \textbf{Role} & \textbf{Lifetime} & \textbf{Primary risk} \\
\midrule
\code{Memory}              & Self-hosted mem0 interface  & Process      & Misconfiguration \\
\code{MemoryClient}        & Cloud mem0 interface        & Process      & Vendor boundary drift \\
\code{MemoryType}          & Write/read taxonomy enum    & Per record   & Wrong policy branch \\
\code{WriteDecision}       & Five-check gate result      & Per write    & Silent bad write \\
\code{RetrievedMemory}     & Retrieval pipeline item     & Per read     & Injection noise \\
\code{WriteAuditRecord}    & Write event payload         & Log window   & Missing provenance \\
\code{RetrievalAuditRecord}& Read event payload          & Log window   & Unexplainable recall \\
\code{ExpiryAuditRecord}   & Deletion event payload      & Log window   & Compliance blind spot \\
\midrule
Episodic memory   & \multicolumn{3}{l}{Session-grounded event, usually TTL-bound; primary risk: staleness} \\
Semantic memory   & \multicolumn{3}{l}{Durable extracted fact; primary risk: conflict drift} \\
Procedural memory & \multicolumn{3}{l}{Reusable workflow rule; primary risk: behavioural overreach} \\
\bottomrule
\end{tabular}
\end{table}

The source layout for this chapter is:

\begin{verbatim}
src/
  agent/
    memory/
      mem0_client.py          # §§8.1, 8.2, 8.13
      memory_types.py         # §8.3
      write_gate.py           # §8.4
      memory_write_node.py    # §§8.5, 8.6, 8.11
      memory_read_node.py     # §8.7
      expiry_scheduler.py     # §§8.8, 8.9
      memory_inspector.py     # §8.10
      audit.py                # §8.12
tests/
  memory/
    test_mem0_integration.py  # §8.14
\end{verbatim}

%% ------------------------------------------------------------
\section{Deployment Mode Must Not Change Call Sites}
\label{sec:deployment-mode}
%% ------------------------------------------------------------

The first design problem is deployment-environment contamination. A naive memory
integration starts with one client locally, switches to a managed service in
staging, and gradually accumulates \code{if~cloud: ... else: ...} forks at
every write and search call site. The deployment decision has leaked into
business logic, and the agent graph becomes untestable in isolation because
each environment requires a different code path. The governing principle is
that deployment mode is a construction-time choice; call sites know only
that they hold a memory object.

mem0 provides two entry points: \code{from~mem0~import~Memory} for
self-hosted deployments wired to a local vector store, and
\code{from~mem0~import~MemoryClient} for the cloud-managed service. The
critical architectural fact is that both classes expose the same operational
surface: \code{m.add()}, \code{m.search()}, \code{m.delete()}, and
\code{m.get\_all()}. An agent that writes to the interface, not to a specific
class, can be redeployed behind either without touching a single node.

This chapter standardises on the self-hosted path throughout because the
book's production model gives operators ownership over data placement,
latency budgets, and deletion guarantees. The cloud client is a valid
deployment alternative; it is not a different programming model.

\begin{lstlisting}[language=Python,
  caption={Both deployment modes expose the same write surface},
  label={lst:mem0-interface}]
# src/agent/memory/mem0_client.py
from mem0 import Memory
from mem0 import MemoryClient

def build_self_hosted(config_dict: dict) -> Memory:
    return Memory.from_config(config_dict)

def build_cloud(api_key: str) -> MemoryClient:
    return MemoryClient(api_key=api_key)

def demo_same_surface(m, user_id: str) -> dict:
    # This call is identical regardless of which builder produced m.
    return m.add(
        "User prefers terse release notes.",
        user_id=user_id,
        metadata={"type": "semantic"},
    )
\end{lstlisting}

Two decisions in Listing~\ref{lst:mem0-interface} deserve explanation.
First, \code{demo\_same\_surface} accepts \code{m} without a type annotation
rather than a union. This is deliberate: annotating as \code{Memory |
MemoryClient} would import both classes at the call site, creating the coupling
the function is supposed to eliminate. The interface contract is enforced by
the shared method surface, not by a type annotation. Second, the listing
demonstrates the sameness of the write call rather than contrasting the
two construction signatures, because sameness is the design property the
rest of the chapter depends on.

A stable interface is not useful if every call site constructs its own client
from different settings; the next section fixes configuration at a single point.

%% ------------------------------------------------------------
\section{Configuration Must Not Live in Code}
\label{sec:config-not-code}
%% ------------------------------------------------------------

The naive wiring for long-term memory is to hardcode a hostname, an
embedder name, and a collection name at the site of the first \code{m.add()}
call. That works until the first migration, at which point half the fleet
writes to one collection and half to another, retrieval quality drops, and
nobody can reconstruct which records belong to which embedder generation.
The governing principle, which Chapter~5 established for LLM settings, applies
equally here: the entire backend configuration is a single validated object
assembled once and injected everywhere as a dependency.

mem0's configuration entry point is \code{Memory.from\_config(config\_dict)},
which expects a dict with four top-level keys: \code{"llm"},
\code{"embedder"}, \code{"vector\_store"}, and \code{"graph\_store"}. Missing
any key raises a runtime error whose message does not identify which key is
absent, making a builder function that makes the structure explicit far safer
than constructing the dict inline. The \code{"llm"} key should point to the
LiteLLM provider introduced in Chapter~5 so that mem0's internal extraction
and conflict-resolution calls pass through the same routing, retry, and cost
accounting layer the rest of the agent uses.

\begin{warningbox}
Changing the embedder after records have been written is a migration, not a
configuration tweak. A Qdrant collection built with 1\,536-dimensional
embeddings cannot be queried safely after switching to a 3\,072-dimensional
model. Treat the embedder model name as schema: version it, and create a new
collection with a migration job rather than silently changing it in settings.
\end{warningbox}

\begin{lstlisting}[language=Python,
  caption={build\_mem0\_config and singleton get\_memory\_client},
  label={lst:mem0-config}]
# src/agent/memory/mem0_client.py
from functools import lru_cache
from mem0 import Memory

COLLECTION_NAME = "agent_memories"

def build_mem0_config(settings) -> dict:
    return {
        "llm": {
            "provider": "litellm",
            "config": {"model": settings.llm_model},
        },
        "embedder": {
            "provider": "openai",
            "config": {
                "model":   "text-embedding-3-small",
                "api_key": settings.openai_api_key.get_secret_value(),
            },
        },
        "vector_store": {
            "provider": "qdrant",
            "config": {
                "collection_name": COLLECTION_NAME,
                "host":            settings.qdrant_host,
                "port":            settings.qdrant_port,
            },
        },
        "graph_store": {
            "provider": "neo4j",
            "config":   {"url": settings.neo4j_url},
        },
    }

@lru_cache(maxsize=1)
def get_memory_client(settings) -> Memory:
    return Memory.from_config(build_mem0_config(settings))
\end{lstlisting}

Three decisions in Listing~\ref{lst:mem0-config} deserve explanation. First,
\code{COLLECTION\_NAME} is a module-level constant, not a field in settings.
This is a deliberate constraint: the collection name is shared infrastructure
across all users and all deployments; making it configurable per-environment
would allow a misconfiguration to silently route writes to the wrong
collection. Second, \code{get\_memory\_client} is wrapped in
\code{@lru\_cache(maxsize=1)}. The \code{Memory} object owns connections to
Qdrant and Neo4j; constructing it per-request would establish and tear down
those connections on every agent invocation. The cache ensures the client is
constructed once per process even in environments where the calling code is
executed many times. Third, the Neo4j graph store is listed explicitly rather
than defaulted to \code{"in\_memory"}. The graph store is what enables mem0's
conflict resolution for semantic memories; an in-memory graph store silently
loses all relationship data on process restart, which is appropriate for tests
but not for any deployment where semantic deduplication matters.

With a singleton client in place and backend choices committed, the next
question is what kind of thing is being written --- because type drives every
downstream policy decision.

%% ------------------------------------------------------------
\section{Memory Type Must Be Explicit Before Every Write}
\label{sec:type-explicit}
%% ------------------------------------------------------------

The temptation is to write everything with a single generic label and
differentiate later in retrieval. This is the same mistake as storing all
events in one untyped log and hoping queries will filter correctly.
When the type is absent or inconsistent at write time, retrieval policy
branches become meaningless, TTL assignment produces wrong lifetimes, and
the injection formatter (Chapter~7 \S\ref{sec:compact-format}) cannot
produce a labelled line because the label is unknown. The governing principle
is that type must be resolved before \code{m.add()} is called, not after.

Chapter~7 introduced \code{MemoryType} as a \code{Literal} string for the
format layer. This section promotes it to a Python \code{Enum} so that the
write gate, the read pipeline, and the audit records can branch on it with
type safety. Three values cover the space of what an agent should persist.
\emph{Episodic} memory is a specific past event grounded in a session: what
the user asked, what the agent decided, what outcome was produced. It is
time-indexed and usually TTL-bounded. \emph{Semantic} memory is a durable
extracted fact about the user or world: ``user's timezone is UTC+9'',
``user prefers imperative prose''. It is not session-indexed; it accumulates
and is updated by mem0's internal conflict resolution. \emph{Procedural}
memory is a workflow rule or learned preference for how to accomplish a
class of tasks: ``when writing SQL, always qualify table names''. It
instructs behaviour rather than recalling facts.

\begin{notebox}
Procedural memories are not injected through Chapter~7's memory block format.
They describe how the agent should behave, not what it should recall; they
belong in the \code{ROLE} slot of Chapter~6's context assembly node, not in
the \code{HISTORY} slot. The injection node
(\S\ref{sec:injection-pipeline} of Chapter~7) skips records whose type
is \code{PROCEDURAL}. The context assembly node (Chapter~6
\S\ref{sec:context-assembly-node}) reads them separately.
\end{notebox}

\begin{lstlisting}[language=Python,
  caption={MemoryType enum and LLM-backed classifier with safe fallback},
  label={lst:memory-types}]
# src/agent/memory/memory_types.py
from enum import Enum
from langchain_core.runnables import Runnable

class MemoryType(str, Enum):
    EPISODIC    = "episodic"
    SEMANTIC    = "semantic"
    PROCEDURAL  = "procedural"

_CLASSIFY_PROMPT = (
    "Classify the following text as exactly one of: "
    "episodic, semantic, procedural. "
    "Reply with the single word only.\n\nText: {content}"
)

def classify_memory_type(
    content: str, llm: Runnable
) -> MemoryType:
    prompt = _CLASSIFY_PROMPT.format(content=content[:400])
    raw    = llm.invoke(prompt).content.strip().lower()
    try:
        return MemoryType(raw)
    except ValueError:
        return MemoryType.EPISODIC
\end{lstlisting}

Two decisions in Listing~\ref{lst:memory-types} deserve explanation. First,
\code{MemoryType} inherits from both \code{str} and \code{Enum}, mirroring the
pattern Chapter~6 used for \code{SlotName}. This makes a \code{MemoryType}
value serialise as a plain string in audit records and Langfuse events without
any custom encoder, while remaining a closed enumeration that the type checker
validates at every call site. Second, the fallback on \code{ValueError} returns
\code{EPISODIC} rather than raising. A classification failure should not abort
the write; \code{EPISODIC} is the safest conservative choice because it carries
the shortest TTL, meaning an incorrectly classified record will be expired
rather than living indefinitely as a semantic or procedural entry.

With type established, the write gate can now apply policy checks that are
specific to each type's semantics.

%% ------------------------------------------------------------
\section{Every Write Must Justify Itself}
\label{sec:write-gate}
%% ------------------------------------------------------------

The problem with writing every assistant turn to long-term memory is not
storage cost --- it is retrieval poisoning. Every trivial exchange, duplicate
preference, and fragmented tool observation that enters the store is a record
that will be returned in future similarity searches, ranked alongside genuinely
useful memories, and possibly injected into a prompt where it crowds out a
more relevant fact. The governing principle is that a write must be approved by
a gate that asks five questions, in order, before any \code{m.add()} call is
made.

The five checks are designed to fail fast: the gate short-circuits on the
first failed check and returns a \code{WriteDecision} with \code{approved=False}
and the name of the check that blocked the write. This makes the audit record
(Section~\ref{sec:audit-trail}) actionable --- an operator can query
``which check blocked the most writes this week'' and tune thresholds
accordingly. The five checks are:

\begin{enumerate}
    \item \textbf{Significance} (\code{is\_significant}): does the content
    exceed a minimum-information threshold? Content shorter than 20 words that
    contains no named entity, no quantitative fact, and no stated preference
    is considered trivial.

    \item \textbf{Type classification} (\code{classify\_memory\_type}): can
    the content be assigned a valid \code{MemoryType}? A classification that
    returns \code{None} on a second pass aborts the write; this should not
    happen with the fallback in \S\ref{sec:type-explicit}, but the gate
    defends against future classifier changes.

    \item \textbf{PII sensitivity} (\code{contains\_pii}): does the content
    match patterns for credit card numbers, social security numbers, raw email
    addresses, or passwords? Any match aborts the write and logs a
    \code{PIIBlockedEvent}. The check uses a regex battery; it is fast and
    imperfect. Do not use it as the sole PII control in a regulated
    environment.

    \item \textbf{Deduplication} (\code{is\_duplicate}): does
    \code{m.search(content, filters=\{"user\_id": uid\}, limit=1)} return a
    result with cosine similarity above 0.92? If so, the content is already in
    the store and the write is redundant.

    \item \textbf{TTL assignment} (\code{assign\_ttl}): what expiry timestamp
    should this record carry? \code{EPISODIC} records expire in 90 days,
    \code{SEMANTIC} in 365 days, and \code{PROCEDURAL} records do not expire.
    The assignment always succeeds; the check exists to make the TTL decision
    explicit and testable.
\end{enumerate}

\begin{lstlisting}[language=Python,
  caption={WriteDecision dataclass and run\_write\_gate short-circuit logic},
  label={lst:write-gate}]
# src/agent/memory/write_gate.py
import re
from dataclasses import dataclass, field
from datetime import datetime, timedelta
from src.agent.memory.memory_types import MemoryType

@dataclass
class WriteDecision:
    approved:        bool
    reason:          str
    blocked_check:   str
    memory_type:     MemoryType | None
    ttl_expires_at:  datetime | None
    importance_score: float = 0.0
    expected_version: int   = 1

_PII_RE = re.compile(
    r"\b\d{4}[- ]?\d{4}[- ]?\d{4}[- ]?\d{4}\b"  # credit card
    r"|\b\d{3}-\d{2}-\d{4}\b"                      # SSN
    r"|password\s*[:=]\s*\S+",
    re.IGNORECASE,
)

def is_significant(content: str) -> bool:
    return len(content.split()) >= 20

def contains_pii(content: str) -> bool:
    return bool(_PII_RE.search(content))

def is_duplicate(content: str, user_id: str, m) -> bool:
    results = m.search(content, filters={"user_id": user_id}, limit=1)
    return bool(results and results[0].get("score", 0.0) >= 0.92)

def assign_ttl(mt: MemoryType, now: datetime) -> datetime | None:
    if mt == MemoryType.EPISODIC:  return now + timedelta(days=90)
    if mt == MemoryType.SEMANTIC:  return now + timedelta(days=365)
    return None  # PROCEDURAL never expires
\end{lstlisting}

\begin{lstlisting}[language=Python,
  caption={run\_write\_gate composing the five checks into a single decision},
  label={lst:write-gate-run}]
# src/agent/memory/write_gate.py (continued)
def run_write_gate(
    content:  str,
    user_id:  str,
    m,
    llm,
    now:      datetime,
) -> WriteDecision:
    if not is_significant(content):
        return WriteDecision(False, "content too short", "significance",
                             None, None)
    mt = classify_memory_type(content, llm)
    if mt is None:
        return WriteDecision(False, "unclassifiable", "type_classification",
                             None, None)
    if contains_pii(content):
        return WriteDecision(False, "PII detected", "pii_sensitivity",
                             mt, None)
    if is_duplicate(content, user_id, m):
        return WriteDecision(False, "duplicate", "deduplication", mt, None)
    ttl = assign_ttl(mt, now)
    return WriteDecision(True, "approved", "", mt, ttl)
\end{lstlisting}

Three decisions in Listings~\ref{lst:write-gate} and~\ref{lst:write-gate-run}
deserve explanation. First, \code{WriteDecision} is a \code{dataclass} rather
than a plain tuple or dict. The downstream write node, the audit emitter, and
the test suite all read named fields; a dataclass makes that access type-safe
and makes the decision's full structure visible in a single class definition.
Second, \code{is\_duplicate} calls \code{m.search()} rather than maintaining a
local cache. The gate runs before the write; the most authoritative answer to
``does this exist'' is always the vector store, not an in-process approximation
that may lag if another session wrote to the same user's memory concurrently.
Third, the PII regex is intentionally coarse. It is the first line of a
defence-in-depth stack, not the whole stack. Organisations subject to GDPR or
HIPAA should layer a dedicated PII classification model on top of this regex
check; the gate's contract is that it blocks obvious patterns fast, not that
it catches everything.

With a gate deciding what can be written, the next question is what metadata
must accompany every approved write so the record carries its own lifecycle
policy.

%% ------------------------------------------------------------
\section{Episodic Writes Must Carry Their Own Lifecycle}
\label{sec:episodic-writes}
%% ------------------------------------------------------------

The problem with a bare \code{m.add(content)} call is that the record has no
lifecycle. It cannot be expired, it cannot be pruned by importance, and it
cannot be retrieved with a reliable timestamp because mem0's internal
\code{created\_at} field may not reflect when in the session the fact was
relevant. The governing principle is that every approved episodic write
carries metadata that makes the record self-describing: what it is, how
important it is, when it should expire, and what version it represents.

The \code{m.add()} signature for a full episodic write passes the conversation
turns as a list of role-dicts rather than a flat string. mem0's internal
extraction pipeline uses the role structure --- distinguishing user statements
from assistant inferences --- to decide what is worth extracting as a semantic
fact. Passing a flat string collapses that distinction and reduces extraction
quality.

\begin{lstlisting}[language=Python,
  caption={write\_episodic\_memory with full metadata and memory\_id capture},
  label={lst:episodic-write}]
# src/agent/memory/memory_write_node.py
from datetime import datetime

def score_importance(
    messages: list[dict],
    query_terms: list[str],
) -> float:
    text   = " ".join(m.get("content","") for m in messages)
    words  = text.split()
    length = min(len(words) / 100.0, 0.4)
    term_hits = sum(1 for t in query_terms if t.lower() in text.lower())
    overlap = min(term_hits / max(len(query_terms), 1), 0.6)
    return round(length + overlap, 3)

def write_episodic_memory(
    messages:   list[dict],
    user_id:    str,
    session_id: str,
    m,
    gate,
    now:        datetime,
    query_terms: list[str] = [],
) -> str | None:
    if not gate.approved:
        return None
    importance = score_importance(messages, query_terms)
    result = m.add(
        messages,
        user_id=user_id,
        metadata={
            "type":           gate.memory_type.value,
            "ttl_expires_at": gate.ttl_expires_at.isoformat()
                              if gate.ttl_expires_at else None,
            "importance_score": importance,
            "session_id":     session_id,
            "version":        gate.expected_version,
        },
    )
    return result.get("memory_id") or result.get("id")
\end{lstlisting}

Two decisions in Listing~\ref{lst:episodic-write} deserve explanation. First,
\code{score\_importance} uses word count and query-term overlap rather than
an LLM call. An LLM-based importance scorer would add latency and cost to
every approved write; a lightweight heuristic that returns a float in
\([0, 1]\) is sufficient for the pruning threshold in \S\ref{sec:pruning}
and for the audit record in \S\ref{sec:audit-trail}. Second,
\code{ttl\_expires\_at} is serialised as an ISO 8601 string rather than stored
as a \code{datetime} object. mem0's metadata payload is ultimately serialised
to Qdrant's JSON payload store; a \code{datetime} object would cause a
serialisation error at write time. The expiry job (\S\ref{sec:ttl-expiry})
parses this string back to a \code{datetime} when it reads records.

The write node persists episodic content explicitly. mem0 also runs its own
extraction pipeline silently after every \code{m.add()}, which the next
section explains and shows how to observe.

%% ------------------------------------------------------------
\section{Semantic Extraction Is Automatic, But Must Be Monitored}
\label{sec:semantic-extraction}
%% ------------------------------------------------------------

The problem is that automatic extraction is invisible unless deliberately
surfaced. After every \code{m.add()} call, mem0 runs an LLM-backed pipeline
that parses the content for durable facts, writes them as semantic records to
the Neo4j graph layer, and handles conflicts against existing semantic records
by updating or deleting nodes rather than duplicating them. If this pipeline
produces empty results on turns that clearly contain factual statements, or
if it creates spurious facts from ambiguous phrasing, the operator has no way
to detect the problem without explicitly logging the extraction output.
The governing principle is that automatic behaviour must be exposed as a
first-class observable event.

mem0 surfaces extraction results in the return value of \code{m.add()} under
the \code{"extracted\_memories"} key, as a list of dicts. Each dict has a
\code{"memory"} field with the extracted text and an \code{"event"} field
with one of four values: \code{"ADD"} (new record created), \code{"UPDATE"}
(existing record modified by conflict resolution), \code{"DELETE"} (existing
record removed as superseded), or \code{"NONE"} (no extraction action taken).
Logging this list to Langfuse converts automatic extraction into auditable
events.

\begin{warningbox}
The PII check at the write gate runs on the raw \code{content} argument before
\code{m.add()} is called. It does not run on extracted semantic records,
which may surface PII fragments that were embedded in longer text the gate
considered clean. Call \code{pii\_scan\_extracted} on every extraction result
before logging it to Langfuse or storing it in audit records.
\end{warningbox}

\begin{lstlisting}[language=Python,
  caption={log\_semantic\_extractions with PII scan and Langfuse event emission},
  label={lst:semantic-extraction}]
# src/agent/memory/memory_write_node.py (continued)
import re

_PII_RE_SIMPLE = re.compile(
    r"\b\d{4}[- ]?\d{4}[- ]?\d{4}[- ]?\d{4}\b|\b\d{3}-\d{2}-\d{4}\b",
)

def pii_scan_extracted(extracted: list[dict]) -> list[dict]:
    return [
        e for e in extracted
        if not _PII_RE_SIMPLE.search(e.get("memory", ""))
    ]

def log_semantic_extractions(
    add_result:      dict,
    langfuse_client,
    user_id:         str,
    trace_id:        str,
) -> list[str]:
    raw = add_result.get("extracted_memories", [])
    clean = pii_scan_extracted(raw)
    events = [e.get("event", "NONE") for e in clean]
    langfuse_client.event(
        trace_id=trace_id,
        name="semantic_extraction",
        input={"user_id": user_id},
        output={"count": len(clean), "events": events,
                "memories": [e.get("memory","") for e in clean]},
    )
    return events
\end{lstlisting}

Two decisions in Listing~\ref{lst:semantic-extraction} deserve explanation.
First, \code{pii\_scan\_extracted} filters rather than redacts. Replacing PII
with a placeholder string (``\texttt{[REDACTED]}'') would create an audit record
that looks clean but carries a semantically broken sentence, which could be
injected into a future prompt as misinformation. Dropping the record entirely
is safer. Second, \code{log\_semantic\_extractions} returns the list of event
strings. The caller stores this list in \code{AgentState["meta"]
["last\_extraction\_events"]}, where the \code{MemoryInspector}
(\S\ref{sec:memory-inspection}) includes it in compliance dumps as evidence
of what the agent inferred from the session.

With writes governed and observed, the chapter now turns to retrieval --- the
part that determines what the model ultimately sees.

%% ------------------------------------------------------------
\section{Retrieval Must Narrow Before It Injects}
\label{sec:retrieval-narrow}
%% ------------------------------------------------------------

The problem with passing mem0's raw search results directly to the context
assembly node is that \code{m.search()} optimises for recall, not for prompt
economy. Setting \code{limit=5} returns the five nearest vectors by cosine
distance, but cosine distance ignores recency, penalises no one for age,
and weights a record from three years ago identically to one from yesterday
if their embeddings are close. The context window costs money; injecting
five records when two would suffice is not caution, it is waste. The
governing principle is that retrieval must reduce a wide candidate set through
explicit, auditable steps before any record reaches a prompt.

Chapter~7 \S\ref{sec:injection-pipeline} established the six-step injection
pipeline as the contract: \code{retrieve\_raw\_memories}, \code{filter\_by\_scope},
\code{score\_and\_rank}, \code{select\_top\_k}, \code{format\_memory\_block},
and \code{inject\_memory\_block}. This section replaces the placeholder stub in
\code{retrieve\_raw\_memories} with a real \code{m.search()} call and adds a
recency decay reranker between retrieval and the top-k cut. All downstream
steps from Chapter~7 are reused unchanged.

\begin{notebox}
Setting \code{limit=20} in \code{m.search()} intentionally over-fetches
relative to the final top-5 selection. Fetching only 5 would bypass the
reranking step: the 5 nearest vectors by cosine distance are not the same
as the 5 best candidates after recency decay. The reranker needs a wider
candidate pool to do useful work.
\end{notebox}

\begin{lstlisting}[language=Python,
  caption={retrieve\_raw\_memories with m.search() and rerank\_with\_recency},
  label={lst:retrieval}]
# src/agent/memory/memory_read_node.py
from datetime import datetime, timezone
from src.agent.memory.mem0_client import safe_search
from src.agent.memory.memory_injection_node import RetrievedMemory

MEM0_SEARCH_LIMIT = 20

def decay(days_old: float, half_life: float = 30.0) -> float:
    return 0.5 ** (days_old / half_life)

def rerank_with_recency(
    memories:        list[RetrievedMemory],
    now:             datetime,
    half_life_days:  float = 30.0,
) -> list[RetrievedMemory]:
    for m in memories:
        created = m["created_at"]
        if created.tzinfo is None:
            created = created.replace(tzinfo=timezone.utc)
        age = (now.replace(tzinfo=timezone.utc) - created).days
        m["score"] = m["score"] * decay(age, half_life_days)
    return sorted(memories, key=lambda x: x["score"], reverse=True)

def retrieve_raw_memories(state, m, user_id: str) -> list[RetrievedMemory]:
    query   = _last_query(state)
    results = safe_search(m, query, user_id, limit=MEM0_SEARCH_LIMIT)
    memories = [
        RetrievedMemory(
            id=r["id"],
            type=r.get("metadata", {}).get("type", "episodic"),
            created_at=_parse_dt(r.get("created_at", "")),
            content=r["memory"],
            score=r.get("score", 0.5),
        )
        for r in results
    ]
    if state.get("history"):
        memories.append(RetrievedMemory(
            id="history-summary", type="episodic",
            created_at=datetime.utcnow(),
            content=state["history"], score=1.0,
        ))
    return memories
\end{lstlisting}

Three decisions in Listing~\ref{lst:retrieval} deserve explanation. First,
\code{decay} applies exponential decay with a 30-day half-life. This means a
record that is 30 days old has its cosine score halved; at 90 days, the score
is one-eighth of its original value. The half-life is a parameter rather than
a constant so that different memory types can use different decay rates
--- episodic memories may warrant a shorter half-life than semantic facts about
a user's long-standing preferences. Second, \code{m["score"]} is mutated
in-place inside \code{rerank\_with\_recency} rather than building a new list.
The \code{RetrievedMemory} TypedDict is a local data-flow object, not shared
state; mutation is acceptable and avoids an extra allocation per candidate.
Third, the history summary from \code{state["history"]} is appended with
\code{score=1.0}. This is intentional: the compressed summary produced by
Chapter~7's summarisation node represents content the agent has already decided
is worth preserving; it should not compete with mem0 records in the recency
ranking.

With retrieval governed, the next operational concern is what happens to
records that have lived long enough to be wrong, misleading, or simply
irrelevant.

%% ------------------------------------------------------------
\section{Expiry Must Run on a Schedule, Not on Demand}
\label{sec:ttl-expiry}
%% ------------------------------------------------------------

The problem with lazy expiry --- checking \code{ttl\_expires\_at} during
retrieval and silently skipping expired records --- is that expired records
remain in the vector store. They occupy index space, they appear in
\code{m.get\_all()} responses, and they can still be returned by
\code{m.search()} in the instant before the filter is applied. More
importantly, lazy expiry is invisible: nobody knows how many expired records
exist, and a compliance request for ``all memories we hold for user X''
includes them. The governing principle is that expiry is a scheduled
maintenance operation that physically removes records, not a retrieval
filter that hides them.

APScheduler's \code{BackgroundScheduler} is a daemon-thread scheduler that
runs jobs inside the application process without requiring a separate worker
or queue service. The scheduler is constructed once at application startup,
the expiry job is registered with a \code{CronTrigger}, and the scheduler is
shut down gracefully on process exit via \code{atexit}. A nightly run at
02:00 UTC avoids peak session traffic in most timezones and keeps the
delete-per-day count small, so Qdrant's write-ahead log does not accumulate
a large delete batch.

\begin{lstlisting}[language=Python,
  caption={expire\_ttl\_memories and build\_expiry\_scheduler with CronTrigger},
  label={lst:ttl-expiry}]
# src/agent/memory/expiry_scheduler.py
import atexit
from datetime import datetime, timezone
from typing import Callable
from apscheduler.schedulers.background import BackgroundScheduler
from apscheduler.triggers.cron import CronTrigger

def expire_ttl_memories(
    m,
    user_ids:  set[str],
    now:       datetime,
) -> int:
    deleted = 0
    for uid in user_ids:
        for record in m.get_all(user_id=uid):
            ttl_str = record.get("metadata", {}).get("ttl_expires_at")
            if not ttl_str:
                continue
            ttl = datetime.fromisoformat(ttl_str).replace(tzinfo=timezone.utc)
            if now.replace(tzinfo=timezone.utc) >= ttl:
                m.delete(record["id"])
                deleted += 1
    return deleted

def build_expiry_scheduler(
    m:             object,
    get_user_ids:  Callable[[], set[str]],
) -> BackgroundScheduler:
    scheduler = BackgroundScheduler()
    scheduler.add_job(
        func    = lambda: expire_ttl_memories(m, get_user_ids(), datetime.utcnow()),
        trigger = CronTrigger(hour=2, minute=0),
        id      = "ttl_expiry",
        replace_existing = True,
    )
    scheduler.start()
    atexit.register(scheduler.shutdown)
    return scheduler
\end{lstlisting}

Three decisions in Listing~\ref{lst:ttl-expiry} deserve explanation. First,
\code{get\_user\_ids} is passed as a callable rather than as a set, because the
set of active users changes over time and the job runs nightly; a stale
snapshot of user IDs would miss new users and may include deleted ones. The
callable queries the current active-user registry (a Redis set or a Postgres
table that the write node updates on every \code{m.add()} call) at job
execution time, not at scheduler construction time. Second, \code{atexit.register(
scheduler.shutdown)} ensures that APScheduler's background thread is joined
cleanly when the process exits, rather than being killed mid-iteration of an
active deletion loop. Without this, a deployment rollover can leave a partial
deletion run with no log entry for the incomplete deletes. Third,
\code{expire\_ttl\_memories} returns the count of deleted records. The caller
constructs an \code{ExpiryAuditRecord} from this count and emits it to
Langfuse; the count is the only field that cannot be derived from the audit
record's other fields, so it must be returned explicitly rather than assumed.

Records that are not yet expired but are never retrieved may still be worth
removing; the next section handles that case.

%% ------------------------------------------------------------
\section{Unread Memories Should Be Archived Before Deletion}
\label{sec:pruning}
%% ------------------------------------------------------------

TTL expiry removes memories because they are old. Importance-threshold pruning
removes memories because they were never useful --- their \code{importance\_score}
fell below a threshold and no retrieval audit record has ever included their ID.
The distinction matters because an expired memory might still be wanted for
audit purposes, whereas a never-retrieved, low-importance memory is just noise.
The governing principle is that pruning is two-phase: archive first, delete
second, so the record is recoverable during the archival window if a compliance
or debugging need surfaces after deletion.

The archive target is S3 rather than a second database table because S3
provides lifecycle policies, bucket-level access controls, and cost-tiered
storage that make it suited for long-term archival data. The archive job writes
a JSON object per deleted record and only calls \code{m.delete()} after
confirming the S3 write succeeded. If the S3 write fails, the record stays in
Qdrant; the next weekly run retries the pair.

\begin{lstlisting}[language=Python,
  caption={archive\_to\_s3, prune\_low\_importance, and add\_pruning\_job},
  label={lst:pruning}]
# src/agent/memory/expiry_scheduler.py (continued)
import json
from dataclasses import dataclass, asdict

@dataclass
class PruneArchiveRecord:
    memory_id:       str
    user_id:         str
    content:         str
    metadata:        dict
    archived_at:     str
    importance_score: float

def archive_to_s3(record: dict, s3_client, bucket: str,
                  user_id: str) -> bool:
    body = json.dumps(asdict(PruneArchiveRecord(
        memory_id=record["id"], user_id=user_id,
        content=record.get("memory",""),
        metadata=record.get("metadata",{}),
        archived_at=datetime.utcnow().isoformat(),
        importance_score=record.get("metadata",{}).get("importance_score",0.0),
    )))
    try:
        s3_client.put_object(
            Bucket=bucket,
            Key=f"archives/{user_id}/{record['id']}.json",
            Body=body,
        )
        return True
    except Exception:
        return False

def prune_low_importance(
    m, s3_client, bucket: str,
    user_ids: set[str], threshold: float = 0.2,
) -> int:
    pruned = 0
    for uid in user_ids:
        for rec in m.get_all(user_id=uid):
            score = rec.get("metadata",{}).get("importance_score", 1.0)
            if score < threshold:
                if archive_to_s3(rec, s3_client, bucket, uid):
                    m.delete(rec["id"])
                    pruned += 1
    return pruned

def add_pruning_job(scheduler, m, s3_client, bucket, get_user_ids):
    scheduler.add_job(
        func=lambda: prune_low_importance(
            m, s3_client, bucket, get_user_ids()),
        trigger=CronTrigger(day_of_week="sun", hour=3),
        id="importance_prune", replace_existing=True,
    )
\end{lstlisting}

Two decisions in Listing~\ref{lst:pruning} deserve explanation. First,
\code{importance\_score} defaults to \code{1.0} when absent from metadata
rather than to \code{0.0}. A missing score means the record was written without
the score field --- either by a legacy path or by a test fixture. Defaulting to
\code{1.0} causes the pruning job to skip such records rather than deleting
them, which is the safer conservative choice. Second, \code{add\_pruning\_job}
accepts the scheduler built by \code{build\_expiry\_scheduler} rather than
constructing its own. Both jobs run on the same scheduler instance, which means
both are started and stopped by the same lifecycle code, and both appear in the
same APScheduler job list for monitoring.

The records in Qdrant and S3 describe what the agent has stored. The next
concern is how to surface that state for compliance and debugging without
routing the retrieval through the agent's production hot path.

%% ------------------------------------------------------------
\section{Inspection Is a Compliance Surface, Not a Debug Tool}
\label{sec:memory-inspection}
%% ------------------------------------------------------------

The problem with using \code{m.get\_all()} directly in ad-hoc scripts is that
it has no access control, no audit trail, no serialisation contract, and no
way to attach the result to a Langfuse trace. In regulated environments, a
``show me everything stored for this user'' query is not debugging --- it is a
GDPR Article~15 right-of-access response, and it must be signed, timestamped,
and auditable. The governing principle is that memory inspection is a
first-class operation implemented as a typed class, not a convenience function.

\code{m.get\_all(user\_id=uid)} is expensive: it returns the full record set
for a user and must not be called inside the agent graph's per-request path.
The three legitimate use cases are a compliance export triggered by an operator
API, a debug event attached to a Langfuse trace when \code{LLMSettings.debug}
is \code{True}, and an evaluation fixture call in the test suite
(\S\ref{sec:test-suite}).

\begin{lstlisting}[language=Python,
  caption={MemoryInspector with compliance export and debug event emission},
  label={lst:memory-inspector}]
# src/agent/memory/memory_inspector.py
import hashlib, hmac, json
from dataclasses import dataclass

@dataclass
class MemoryExportRecord:
    memory_id:       str
    content:         str
    type:            str
    created_at:      str
    ttl_expires_at:  str | None
    importance_score: float

class MemoryInspector:
    def __init__(self, m, langfuse_client) -> None:
        self._m  = m
        self._lf = langfuse_client

    def get_all_for_user(self, user_id: str) -> list[dict]:
        return self._m.get_all(user_id=user_id)

    def emit_debug_event(self, user_id: str, state: dict) -> None:
        records = self.get_all_for_user(user_id)
        self._lf.event(
            trace_id = state["meta"].get("langfuse_trace_id"),
            name     = "memory_inspector",
            input    = {"user_id": user_id},
            output   = {
                "count":             len(records),
                "records":           records,
                "extraction_events": state["meta"].get(
                    "last_extraction_events", []),
            },
        )

    def export_for_gdpr(
        self, user_id: str, signing_key: bytes,
    ) -> tuple[list[MemoryExportRecord], str]:
        records = self.get_all_for_user(user_id)
        exports = [MemoryExportRecord(
            memory_id=r["id"],
            content=r.get("memory",""),
            type=r.get("metadata",{}).get("type",""),
            created_at=r.get("created_at",""),
            ttl_expires_at=r.get("metadata",{}).get("ttl_expires_at"),
            importance_score=r.get("metadata",{}).get("importance_score",0.0),
        ) for r in records]
        payload = json.dumps([e.__dict__ for e in exports], sort_keys=True)
        sig = hmac.new(signing_key, payload.encode(), hashlib.sha256).hexdigest()
        return exports, sig
\end{lstlisting}

Two decisions in Listing~\ref{lst:memory-inspector} deserve explanation. First,
the GDPR export signs the JSON payload with an HMAC rather than with a
public-key signature. HMAC is sufficient for integrity verification between the
operator and the user when both parties share the signing key; it avoids the
certificate infrastructure that asymmetric signing would require. The user
receives the payload and the hex signature; they can verify with any SHA-256
HMAC implementation. Second, \code{emit\_debug\_event} reads the trace ID from
\code{state["meta"]["langfuse\_trace\_id"]} rather than constructing a new
trace. The debug event must be attached to the session trace that produced the
memory state being inspected, so the Langfuse timeline shows it in context
with the writes and retrievals that occurred during the same session.

Inspection reveals what exists. The next section addresses a problem that
inspection can surface but cannot fix: two sessions writing conflicting updates
to the same record.

%% ------------------------------------------------------------
\section{Concurrent Writes to the Same Record Need a Version Check}
\label{sec:optimistic-locking}
%% ------------------------------------------------------------

The concurrency problem is subtle because mem0 does not expose transactions.
When two sessions for the same user run simultaneously --- which is common in
multi-tab browser clients or background polling agents --- both may retrieve the
same semantic record, both may modify it based on what they see, and both may
call \code{m.add()} on the updated content. mem0's deduplication check is not
an atomic compare-and-swap; it is a similarity search. The later write wins
and the earlier write's update is lost, silently. The governing principle is
that the record must carry its own version counter, and a write must abort if
the version it expected to overwrite has already changed.

Optimistic locking is the appropriate choice here rather than pessimistic
locking (holding a distributed lock around the read-modify-write cycle)
because memory writes are infrequent relative to read-only operations, the
cost of a retry is low, and a distributed lock would add a Redis or Zookeeper
dependency that is disproportionate for a background write. The version field
is stored in \code{metadata["version"]} and incremented on every successful
write. The write node reads the current version at the start of the write
attempt, embeds the incremented value in the write, and aborts if a
concurrent write has already incremented the version.

\begin{lstlisting}[language=Python,
  caption={versioned\_write with optimistic-lock retry loop},
  label={lst:optimistic-lock}]
# src/agent/memory/memory_write_node.py (continued)

class MemoryConflictError(RuntimeError):
    pass

def fetch_current_version(
    memory_id: str, user_id: str, m
) -> int:
    results = m.search(memory_id, filters={"user_id": user_id}, limit=1)
    if not results:
        return 0
    return results[0].get("metadata", {}).get("version", 0)

def versioned_write(
    messages:          list[dict],
    user_id:           str,
    m,
    gate,
    expected_version:  int,
    max_retries:       int = 3,
) -> str | None:
    if not gate.approved:
        return None
    for attempt in range(max_retries):
        current = fetch_current_version(
            messages[0].get("content","")[:80], user_id, m
        )
        if current != expected_version:
            expected_version = current
            continue
        result = m.add(
            messages, user_id=user_id,
            metadata={
                "type":            gate.memory_type.value,
                "ttl_expires_at":  gate.ttl_expires_at.isoformat()
                                   if gate.ttl_expires_at else None,
                "version":         expected_version + 1,
            },
        )
        return result.get("memory_id") or result.get("id")
    raise MemoryConflictError(
        f"Write conflict unresolved after {max_retries} retries "
        f"for user {user_id}"
    )
\end{lstlisting}

Two decisions in Listing~\ref{lst:optimistic-lock} deserve explanation. First,
the retry loop updates \code{expected\_version} to the fetched \code{current}
value on a mismatch rather than keeping the original expected value. This
means the second retry reads a fresh version rather than re-attempting a
write against a version that is already known to be stale. The content being
written is the same; only the version metadata changes. Second,
\code{MemoryConflictError} is a \code{RuntimeError} rather than a
\code{ValueError}. A conflict after three retries is not a validation
error --- the content is valid. It is a runtime condition that signals
unusual write contention for this user; the calling node should log it and
skip the write rather than propagating the exception up the graph stack and
terminating the session.

With write correctness addressed, the chapter now turns to audit: making every
memory operation visible for debugging, compliance, and evaluation.

%% ------------------------------------------------------------
\section{Every Memory Operation Produces a Structured Event}
\label{sec:audit-trail}
%% ------------------------------------------------------------

The problem with relying on unstructured log lines for memory observability is
that log lines are parse-dependent. Answering ``how many writes did the PII
check block this week for user X?'' requires writing a log parser that is
specific to a log format that may change. A structured Langfuse event with
typed fields can be queried directly, aggregated across sessions, and joined
with retrieval events to build attribution chains. The governing principle is
that every write, retrieval, and expiry operation emits a structured event with
fields that support the four questions operators actually ask: what happened,
to whose data, when, and why.

The three audit record types are \code{WriteAuditRecord},
\code{RetrievalAuditRecord}, and \code{ExpiryAuditRecord}. They are
\code{dataclass} objects rather than plain dicts so that the emitter function
can pass them as typed arguments and the test suite can assert on named fields.

\begin{lstlisting}[language=Python,
  caption={Three audit dataclasses and their Langfuse emitters},
  label={lst:audit-records}]
# src/agent/memory/audit.py
from dataclasses import dataclass
from datetime import datetime
from typing import Literal

@dataclass
class WriteAuditRecord:
    memory_id:        str | None
    user_id:          str
    session_id:       str
    memory_type:      str
    approved:         bool
    blocked_check:    str | None
    importance_score: float
    ttl_expires_at:   str | None
    conflict_retries: int
    timestamp:        datetime

@dataclass
class RetrievalAuditRecord:
    user_id:                  str
    session_id:               str
    query:                    str
    candidate_count:          int
    selected_ids:             list[str]
    injected:                 bool
    injection_blocked_reason: str | None
    timestamp:                datetime

@dataclass
class ExpiryAuditRecord:
    memory_id:        str
    user_id:          str
    expiry_type:      Literal["ttl", "importance"]
    importance_score: float | None
    archived_to_s3:   bool
    timestamp:        datetime
\end{lstlisting}

\begin{lstlisting}[language=Python,
  caption={emit\_ functions attaching audit records to Langfuse traces},
  label={lst:audit-emitters}]
# src/agent/memory/audit.py (continued)
from dataclasses import asdict

def emit_write_audit(record: WriteAuditRecord, lf, trace_id: str) -> None:
    lf.event(trace_id=trace_id, name="memory_write_audit",
             input={"user_id": record.user_id,
                    "approved": record.approved,
                    "blocked_check": record.blocked_check},
             output=asdict(record))

def emit_retrieval_audit(
    record: RetrievalAuditRecord, lf, trace_id: str
) -> None:
    lf.event(trace_id=trace_id, name="memory_retrieval_audit",
             input={"user_id": record.user_id, "query": record.query},
             output=asdict(record))

def emit_expiry_audit(record: ExpiryAuditRecord, lf) -> None:
    # Expiry runs outside any session; no trace_id is available.
    lf.event(name="memory_expiry_audit",
             input={"memory_id": record.memory_id,
                    "expiry_type": record.expiry_type},
             output=asdict(record))
\end{lstlisting}

Two decisions in Listings~\ref{lst:audit-records} and~\ref{lst:audit-emitters}
deserve explanation. First, \code{ExpiryAuditRecord} events are emitted without
a \code{trace\_id}. Expiry runs from a scheduled background job that has no
session context; creating a synthetic trace ID would make the events appear
in the Langfuse UI as orphaned sessions, confusing the session-level analytics.
A trace-ID-free event lands in Langfuse's global event stream, where it can be
queried independently. Second, both the \code{input} and \code{output} fields
of each event carry the full dataclass as the output while projecting only the
most-queried fields into \code{input}. Langfuse's dashboard filter works on
\code{input} fields; surfacing \code{user\_id}, \code{approved}, and
\code{query} at the top level makes it possible to build compliance dashboards
without parsing the full output blob.

The audit trail answers questions about individual operations. The final
system-level question is whether operations from one user can ever touch
another user's data.

%% ------------------------------------------------------------
\section{Tenant Isolation Belongs at Every Call Site}
\label{sec:tenant-isolation}
%% ------------------------------------------------------------

The temptation is to enforce isolation architecturally --- by creating a
separate Qdrant collection per user --- and then trust that the infrastructure
prevents cross-tenant access. This is the wrong pattern for two reasons.
Creating thousands of Qdrant collections for thousands of users incurs per-collection
overhead in the Qdrant service, makes cross-user analytics (model evaluation,
usage reporting) require fan-out over every collection, and encodes user
identifiers into infrastructure resource names, creating a data governance
problem when users are deleted or merged. The governing principle is that
isolation is a single-collection problem solved by enforcing \code{user\_id}
as a mandatory filter at every call site.

\begin{warningbox}
Never call \code{m.search()} without \code{filters=\{"user\_id": uid\}}.
A search without the filter returns records from every user in the collection.
mem0 does not raise an error when the filter is absent; it silently returns
cross-tenant results. The \code{safe\_search} wrapper below makes the
omission a hard error rather than a silent data leak.
\end{warningbox}

\begin{lstlisting}[language=Python,
  caption={safe\_search, safe\_add, and safe\_delete enforcing user\_id at every call},
  label={lst:tenant-isolation}]
# src/agent/memory/mem0_client.py (continued)

USER_ID_REQUIRED_ERROR = "user_id is required for all mem0 operations"

def safe_search(
    m, query: str, user_id: str, limit: int = 20
) -> list[dict]:
    assert user_id, USER_ID_REQUIRED_ERROR
    return m.search(query, filters={"user_id": user_id}, limit=limit)

def safe_add(
    m, messages: list[dict] | str, user_id: str, metadata: dict
) -> dict:
    assert user_id, USER_ID_REQUIRED_ERROR
    metadata.setdefault("user_id", user_id)
    return m.add(messages, user_id=user_id, metadata=metadata)

def safe_delete(m, memory_id: str, user_id: str) -> None:
    assert user_id, USER_ID_REQUIRED_ERROR
    results = m.search(memory_id[:80], filters={"user_id": user_id}, limit=1)
    if not results:
        return
    found_id = results[0].get("id")
    if found_id == memory_id:
        m.delete(memory_id)
\end{lstlisting}

Two decisions in Listing~\ref{lst:tenant-isolation} deserve explanation. First,
\code{safe\_delete} verifies ownership before deleting. mem0's \code{m.delete()}
accepts any memory ID without checking who owns it; passing the wrong ID ---
due to a bug in the ID propagation chain --- would silently delete another
user's record. The ownership check is a similarity search rather than a direct
ID lookup because mem0 does not expose a \code{get\_by\_id} method at this API
level. Second, \code{safe\_add} writes \code{user\_id} into the metadata dict
as a redundant field in addition to the \code{user\_id} keyword argument. This
makes the tenant identifier present in the Qdrant payload directly, where it
can be filtered using Qdrant's native payload filter API for analytics queries
that bypass the mem0 client.

With isolation guaranteed at every call site, the only remaining gap is
demonstrating that the entire stack behaves correctly under adversarial
conditions. The test suite closes that gap.

%% ------------------------------------------------------------
\section{Tests Must Prove Policy, Not Just Coverage}
\label{sec:test-suite}
%% ------------------------------------------------------------

The problem with a test suite that achieves high line coverage but only tests
the happy path is that it proves the gate compiles, not that it governs.
The five write-gate checks, the retrieval reranker, the expiry job, the
version conflict, and the multi-tenant isolation all have failure modes that
only manifest when specific conditions are met: a short message, a high
similarity score, an expired timestamp, a concurrent version bump, or a
missing \code{user\_id}. The governing principle is that each test targets one
failure mode with a fixture constructed to trigger it.

All tests mock \code{Memory} using \code{unittest.mock.MagicMock}. No live
Qdrant or Neo4j instance is required. \code{m.add()} returns
\code{\{"memory\_id": "test-id-1"\}}; \code{m.search()} returns a configurable
fixture list; \code{m.get\_all()} returns a fixed list. The mock isolates the
gate and pipeline logic from network latency and makes fixture construction
deterministic.

\begin{lstlisting}[language=Python,
  caption={mock\_memory fixture and two representative policy tests},
  label={lst:test-suite}]
# tests/memory/test_mem0_integration.py
import pytest
from unittest.mock import MagicMock
from datetime import datetime, timedelta
from src.agent.memory.write_gate import run_write_gate
from src.agent.memory.expiry_scheduler import expire_ttl_memories

@pytest.fixture
def mock_memory():
    m = MagicMock()
    m.add.return_value   = {"memory_id": "test-id-1"}
    m.search.return_value = []
    m.get_all.return_value = []
    m.delete.return_value  = None
    return m

@pytest.fixture
def mock_llm():
    llm = MagicMock()
    llm.invoke.return_value = MagicMock(content="episodic")
    return llm

def test_pii_blocks_credit_card(mock_memory, mock_llm):
    content = ("User's card is 4111 1111 1111 1111 "
               "and they want to pay for the subscription.")
    decision = run_write_gate(content, "u1", mock_memory, mock_llm,
                              datetime.utcnow())
    assert not decision.approved
    assert decision.blocked_check == "pii_sensitivity"
    mock_memory.add.assert_not_called()

def test_ttl_expiry_deletes_only_expired(mock_memory):
    now  = datetime.utcnow()
    past = (now - timedelta(days=1)).isoformat()
    future = (now + timedelta(days=1)).isoformat()
    mock_memory.get_all.return_value = [
        {"id": "old-1", "metadata": {"ttl_expires_at": past}},
        {"id": "new-1", "metadata": {"ttl_expires_at": future}},
        {"id": "no-ttl", "metadata": {}},
    ]
    deleted = expire_ttl_memories(mock_memory, {"u1"}, now)
    assert deleted == 1
    mock_memory.delete.assert_called_once_with("old-1")
\end{lstlisting}

Two decisions in Listing~\ref{lst:test-suite} deserve explanation. First, the
\code{mock\_llm} fixture returns \code{"episodic"} as the classification
response. This is not an arbitrary choice; it ensures the type-classification
check passes and the test can reach the PII check, which is the check under
test. Returning an invalid classification string would cause the test to fail
at the wrong check, making it a test of the classifier fallback rather than of
the PII gate. Second, \code{test\_ttl\_expiry\_deletes\_only\_expired} includes
a record with no \code{ttl\_expires\_at} field alongside the expired and
non-expired records. This tests the guard clause in \code{expire\_ttl\_memories}
that skips records without a TTL --- a guard that exists because procedural
memories are written without a TTL and must never be deleted by the expiry
job.

Chapter~9 extends this test suite into full evaluation: multi-turn session
fixtures, precision@5 measurement against a labelled retrieval dataset, and a
prompt-drift regression test that verifies model output does not shift when
memory injection is disabled versus enabled. The infrastructure built in this
chapter --- the write gate, the retrieval pipeline, the audit records, and the
safe call wrappers --- becomes the fixture surface that Chapter~9's session
replay harness queries to attribute output changes to specific memory events.
